% Part:sets-functions-relations
% Chapter: size-of-sets
% Section: reduction-alt

\documentclass[../../../include/open-logic-section]{subfiles}

\begin{document}

\olfileid[ar]{sfr}{siz}{red-alt}

\olsection{الاختزال}

\begin{editorial}
نثبت في هذا القسم عدم القابلية للتعداد بالاختزال، ونحصل على نتائج \olref[nen-alt]{sec}. وأما العرض البديل الذي يزيد شيئًا من التفصيل ويوافق نتائج \olref[nen]{sec}، ففي \olref[red]{sec}.
\end{editorial}

ثبت بالحجة القطرية أن $\Bin^\omega$ !!{nonenumerable}، وثبت بحجة من جنسها أن $\Pow{\Nat}$ !!{nonenumerable}. غير أن لنا إلى إثبات كون $\Pow{\Nat}$ !!{nonenumerable} سبيلًا آخر: نبين أنه \emph{لو كانت $\Pow{\Nat}$ !!{enumerable} لكانت $\Bin^\omega$ !!{enumerable} أيضًا}. وقد علمنا أن $\Bin^\omega$ !!{nonenumerable}، فيلزم مثل ذلك في $\Pow{\Nat}$.

وهذا \emph{اختزال} مسألة إلى أخرى؛ فقد رددنا تعداد $\Bin^\omega$ إلى تعداد $\Pow{\Nat}$. فحل المسألة الأخيرة، أعني تعداد $\Pow{\Nat}$، لو وجد لأمكن أن نستخرج منه حل الأولى، أعني تعداد $\Bin^\omega$.

ومتى أردنا رد تعداد مجموعة~$B$ إلى تعداد مجموعة~$A$، وجب أن نعين طريقًا يحول تعداد~$A$ إلى تعداد~$B$. وأيسر ذلك أن نضع !!a{surjection} $f\colon A \to B$. فلو كانت $x_1$، $x_2$، \dots{} تعدد~$A$، لكانت $f(x_1)$، $f(x_2)$، \dots{} بعد حذف التكرار وإعادة ترقيم الباقي تعداد~$B$. فمطلبنا هنا !!a{surjection} $f\colon \Pow{\Nat}
\to \Bin^\omega$.

\textbf{تنبيه تصحيحي على الأصل (C040-01).} قد تتكرر قيم الدالة الشاملة؛ فلا تكون
قائمة الصور تعدادًا تقابليًا على التعريف البديل إلا بعد حذف المكرر وإعادة
فهرسة ما بقي.

\begin{prob}
إذا وجدت دالة~!!{injective} $g\colon B \to A$، وكانت $B$~!!{nonenumerable}، فبين أن $A$ كذلك. وذلك بأن تستعمل~$g$ لتحويل تعداد~$A$ المفروض إلى تعداد~$B$.
\end{prob}

\begin{proof}[برهان {\olref[nen-alt]{thm:nonenum-pownat}} بالاختزال]
لنفرض أن $\Pow{\Nat}$ !!{enumerable}، ونأخذ لها تعدادًا $N_{1}$، $N_{2}$، $N_{3}$، \dots

ونضع الدالة $f \colon \Pow{\Nat} \to \Bin^\omega$ بحيث تكون $f(N)$ السلسلة $s$ المعينة بالشرط $s(n) = 1$ إذا وفقط إذا كان $n \in N$، وبالشرط $s(n) = 0$ فيما عدا ذلك.

وهذا تعيين صحيح لدالة؛ إذ متى كان $N \subseteq \Nat$، فكل $n \in \Nat$ إما !!a{element} من~$N$ وإما خارج عنها. فمثلًا مجموعة الأعداد الطبيعية الزوجية $2\Nat = \Setabs{2n}{n \in \Nat} = \{0,2, 4, 6,
\dots\}$ ترسل إلى $1010101\dots$؛ والمجموعة $\emptyset$ ترسل إلى $0000\dots$؛ والمجموعة $\Nat$ ترسل إلى $1111\dots$.

والدالة !!{surjective} أيضًا، لأن كل سلسلة من~$0$ و~$1$ تعين مجموعة الأعداد الطبيعية التي يشغل~$1$ مواضعها، فتكون السلسلة صورة تلك المجموعة. وللتفصيل، إذا كان $s \in \Bin^\omega$ نعين $N
\subseteq \Nat$ بوضع
\[
N = \Setabs{n \in \Nat}{s(n) = 1}
\]
فيكون $f(N) = s$، كما يبينه تعريف~$f$.

ولننظر في القائمة
\[
f(N_1), f(N_2), f(N_3), \dots
\]
إن $f$ !!{surjective}، فكل عنصر من $\Bin^\omega$ قيمة لـ~$f$ عند وسيط ما؛ وهذا الوسيط وارد في التعداد، فترد قيمته في القائمة الجديدة. فهي إذن تستوفي~$\Bin^\omega$ كلها.

فقد تبين أن كون $\Pow{\Nat}$ !!{enumerable} يستلزم كون $\Bin^\omega$ !!{enumerable}. لكن $\Bin^\omega$ !!{nonenumerable} (\olref[nen-alt]{thm:nonenum-bin-omega})، فلزم أن $\Pow{\Nat}$ !!{nonenumerable}.
\end{proof}

\begin{editorial}
\textbf{تنبيه تصحيحي على الأصل (OLSIZ-010).}
سمّى الأصل المتتالية $s_k$ من غير أن يربط $k$؛ فاستُعمل الاسم
$s$ في التعريف كله، إذ تحدد المجموعة $N$ متتاليتها المميزة وحدها.
\end{editorial}

%\begin{explain}
%It is easy to be confused about the direction the reduction goes in.
%For instance, !!a{surjective} function $g \colon \Bin^\omega \to X$
%does \emph{not} establish that $X$ is !!{nonenumerable}.  (Consider $g
%\colon \Bin^\omega \to \Bin$ defined by $g(s) = s(1)$, the function
%that maps a sequence of $0$'s and $1$'s to its first !!{element}.  It
%is surjective, because some sequences start with $0$ and some start
%with $1$. But $\Bin$ is finite.)  Note also that the function $f$ must
%be surjective, or otherwise the argument does not go through:
%$f(x_1)$, $f(x_2)$, \dots{} would then not be guaranteed to include
%all the !!{element}s of~$Y$. For instance, $h\colon \Nat \to
%\Bin^\omega$ defined by
%\[
%h(n) = \underbrace{000\dots0}_{\text{$n$ $0$'s}}
%\]
%is a function, but $\Nat$ is !!{enumerable}.
%\end{explain}

\begin{prob}\label{sfr:siz:red-alt:prob:nat-nat}
أقم برهانًا بالاختزال على أن المجموعة~$X$ لجميع الدوال $f\colon \Nat \to \Nat$ !!{nonenumerable}. وتلميحه إيجاد دالة شاملة من $X$ إلى~$\Bin^\omega$.
\end{prob}

\begin{prob}
بين بالاختزال أن مجموعة جميع \emph{المجموعات المؤلفة من} أزواج الأعداد الطبيعية، أي $\Pow{\Nat \times \Nat}$، !!{nonenumerable}.
\end{prob}

\begin{prob}
أقم حجة اختزال على أن $\Nat^\omega$ !!{nonenumerable}؛ والمراد بها مجموعة جميع المتتاليات اللانهائية من الأعداد الطبيعية.
\end{prob}

%\begin{prob}
%Let $P$ be the set of functions from $\Nat$ to the set $\{0\}$, and let $Q$ be the set of \emph{partial}
%functions from the set of positive integers to the set $\{0\}$. Show
%that $P$~is !!{enumerable} and $Q$~is not. (Hint: reduce the problem
%of enumerating $\Bin^\omega$ to enumerating~$Q$).
%\end{prob}

\begin{prob}
إذا كانت $S$ مجموعة جميع !!{surjection}s من~$\Nat$ إلى $\{0,1\}$، أي كانت $S$ تضم جميع !!{surjection}s~$f \colon \Nat
\to \Bin$، فبرهن أن $S$ !!{nonenumerable}.
\end{prob}

\begin{prob}
بين أن $\Real$، وهي مجموعة الأعداد الحقيقية، !!{nonenumerable}.
\end{prob}

\end{document}
